\subsection{Explicit CHSH violation from non-injective projection.}
  \label{subsec:explicit-chsh-violation-from-non-injective-projection}
  We make the Bell--CHSH violation explicit by exhibiting an observable probability
  model generated by a non-injective projected description.

  Let measurement settings be $x\in\{0,1\}$ at $A$ and $y\in\{0,1\}$ at $B$,
  and outcomes be $a,b\in\{-1,+1\}$.
  Consider the no-signalling correlations defined by
  \begin{equation}
    P(a,b|x,y)
    \;=\;
    \begin{cases}
      \tfrac12 & \text{if } ab = (-1)^{xy},\\
      0 & \text{otherwise.}
    \end{cases}
    \label{eq:prbox_distribution}
  \end{equation}
  Equivalently, for each setting pair $(x,y)$, the outcomes are perfectly
  correlated when $xy=0$ and perfectly anti-correlated when $xy=1$, while each
  marginal is uniform: $P(a|x)=P(b|y)=1/2$.

  The correlators are therefore
  \begin{equation}
    E_{xy} \;\equiv\; \sum_{a,b=\pm 1} ab\,P(a,b|x,y) \;=\; (-1)^{xy},
    \label{eq:prbox_correlators}
  \end{equation}
  and the CHSH expression is
  \begin{equation}
    S \;\equiv\; \bigl|E_{00}+E_{01}+E_{10}-E_{11}\bigr|
    \;=\; |1+1+1-(-1)| \;=\; 4 \;>\; 2.
    \label{eq:chsh_violation}
  \end{equation}
  This violates the Bell--CHSH bound while preserving operational no-signalling.

  \medskip
  \noindent
  \textit{Realization as a projected description.}
  Define an underlying finite configuration space
  $\Omega \equiv \{(a_0,a_1,b_0,b_1)\,|\,a_x,b_y\in\{-1,+1\}\}$ equipped with a
  normalized measure $\mu$ supported only on the subset satisfying the global
  constraint
  \begin{equation}
    a_0 b_0 = a_0 b_1 = a_1 b_0 = +1, \qquad a_1 b_1 = -1.
    \label{eq:global_constraint}
  \end{equation}
  Such a constraint is admissible here because $\Omega$ is not assumed to factorize
  into independent subsystem state spaces at the underlying level
  (Appendix~B.1--B.3).
  Observable pairs are obtained through a context-dependent projection
  \begin{equation}
    \Pi_{x,y}:\Omega\to O_A\times O_B,\qquad
    \Pi_{x,y}(a_0,a_1,b_0,b_1) \;=\; (a_x,b_y).
    \label{eq:contextual_projection}
  \end{equation}
  The dependence of the projection on the measurement settings
  does not introduce contextual hidden variables.
  It merely reflects that different experimental settings select
  different components of the same underlying configuration $\omega$,
  analogously to reading different coordinates of a fixed vector.
  Since many distinct underlying quadruples map to the same observable pair
  $(a,b)$ for fixed $(x,y)$, the map $\Pi_{x,y}$ is generally non-injective.
  The observable probabilities are then exactly of the projected form
  \begin{equation}
    P(a,b|x,y) \;=\; \mu\!\left(\Pi_{x,y}^{-1}(a,b)\right),
    \label{eq:projected_probabilities_chsh}
  \end{equation}
  and, for a symmetric choice of $\mu$ on the constrained subset
  Eq.~\eqref{eq:global_constraint}, reproduce Eq.~\eqref{eq:prbox_distribution}.

  \medskip
  \noindent
  This explicit construction illustrates the logical point emphasized throughout:
  Bell--CHSH violations may arise because the effective observable description is
  generated by a non-injective projection from an underlying configuration space
  that does not admit a subsystem factorization compatible with Bell-type probabilistic assumptions.
  No superluminal influence is invoked at the level of underlying configurations;
  the violation reflects the structural mismatch between factorizable effective
  models and projected descriptions with irreducible global constraints.
